\chapter{Conclusion}\label{ch:conclusion}
\graphicspath{{../Figures/80-conclusion/}{figs/80-conclusion/}}

This dissertation moved the biomimetic humanoid robot project from an unvalidated design concept to a quantitatively characterized actuator set and experimentally tested knee-torque models, and it defined the path from those experimental foundations to neuromechanical control.

At the actuator level, the analysis in Chapter~\ref{ch:muscles} produced novel equations for calculating isometric force in Festo \qtylist{10;20}{\mm}-diameter BPAs. The work elucidated the relationship between maximum BPA force at \qty{620}{\kPa} and resting length, $F_{620}(l_0)$, a previously unreported characteristic of these actuators that the manufacturer's tool predicts with the opposite trend, and it produced a nondimensional force surface, $F^*(\epsilon^*,P^*)$, with only three diameter-dependent coefficients. Taken together, these fits predict isometric BPA force across resting lengths from a small number of measurements, correct the large over-predictions of the manufacturer's tool at short resting lengths, and outperform published single-length models. In a design workflow, these equations make it possible to select actuator diameter and resting length and to evaluate whether a proposed muscle routing can meet force and range-of-motion requirements, before committing to a prototype.

At the joint level (Chapter~\ref{ch:torque}), the BPA model was combined with transmission corrections for effective length, series compliance, and wrapping. Identification on the pinned knee improved torque prediction, and the shared correction terms also improved predictions on the biomimetic four-bar knee. The corrected model supports route redesign: the optimized two-BPA flexor and cam-extensor configurations are predicted to meet their human-scaled torque targets across the design range. The isometric tests demonstrate that the redesigned configurations meet or exceed the human reference torques over most of the tested range of motion.

The sensory-afferent resource of Chapter~\ref{ch:database} makes relevant physiological reports easier to locate and compare while retaining their experimental context. Chapter~\ref{ch:sns} provides a framework for connecting neural controllers to mechanical models and testing them under specified support, stimulation, and perturbation conditions. Independent implementations reproduced the contact-driven rhythm core with exact period agreement and trace correlations above 0.9999. BPA-actuated models demonstrated air stepping and supported standing. The coupled pathway and perturbation tests identified model-specific feedback dependence and posture limitations. These results establish implemented experimental capabilities; harness-free walking remains unachieved.

The next steps are to complete model integration, extend the mechanical characterization to dynamic conditions, improve load-bearing and balance control, and test the resulting system under progressively reduced support (Chapter~\ref{ch:overall}). These steps are required before using whole-body robot behavior as evidence about a biological control hypothesis.

The dissertation provides experimentally grounded actuator and joint models, a curated afferent resource, and simulation tools for that work. Their scientific value lies in making assumptions explicit and testable as the project moves toward controlled neuromechanical experiments on a physical robot.


